This chapter explains how the system was technically designed,
implemented, and improved during the project. Its purpose is not only to
show what was built, but also how the group worked technically and why
specific engineering choices were made.

% --------------------------------
\section{Frontend Implementation}

Frontend for this project was implemented in C\# with Avalonia UI framework and LiveCharts for graph visualization. MVVM pattern and data bindings were used with Avalonia. Frontend is split into supporting elements and pages.\par
Supporting elements include sidebar navigation panel, top navigation tab bar and chart component. Sidebar navigation panel has the logo of our project and a list of scenarios which can be switched by clicking on buttons. Top navigation tab bar contains tab names for each page and a combobox for selecting different periods. The period checkbox is separate for each scenario. The chart component is used for dynamically adding charts to the page. It is used in the optimizer page.\par
There are 3 pages - overview, optimizer and production units page. After our initial Figma prototype, we had 4 pages, the fourth one being the configuration page which handled scenario specific data. However, later down the line we have merged the production units page and configuration page into one production units page because it improved user experience, clarity and reduced clutter. The first page, overview page, contains a map of the Heatington city and various general graphs such as heat demand and production, electricity prices which are populated when the optimizer is run. The second page, optimizer page, has a button to start the optimizer and to export the results of the optimization. Furthermore, after the optimization completion, it shows totals for costs, revenue, amount of heat produced and consumed, energy used and amount of CO2 emissions. Maintenance periods are shown. The graphs with detailed optimizer information such as unit scheduling, price for 1 MWh for each unit at a given hour and costs for each unit at the given hour are shown. It is possible to use a combobox for selecting specific units to be shown in the graphs. In the third production units page, all of the production units are shown with their properties - image, name, heat production, electricity, minimum and maximum maintenance hours, can the unit be maintained, unit colour in the graphs and so on. In the same page, production units can be selected or deselected for the scenario, units can be edited, deleted or new production units can be added.\par
Each component and page have their respective view models responsible for filling the data required for the component. All of the data and actions performed go through Data Manager directly or indirectly. Much of the data displayed in the UI is displayed directly without any modification and a small part of data shown is heavily processed data with inputs coming from Data Manager. Such data include data provided to the charts which contains numerous LINQ queries, handling unit selection, dynamic colors and empty data periods.\par
During our frontend implementation, we have faced major challenges writing clean code with Avalonia UI framework. First of all, data which is stored outside of Avalonia is difficult to update reactively without desynchronization issues. For example, Data Manager provides result data for the optimization and the result is stored outside of Avalonia code because it is part of the domain layer. Avalonia does not update the UI with new data automatically. Instead, C\# events could be passed around through the entire hierarchy of UI components to the Optimizer view model to call a refresh method once new data is available. However, such “prop drilling” to pass a refresh event was proven difficult as getting data to the view model from other components would have required changing every view and view model along the UI hierarchy. Sadly, some less efficient reactivity methods needed to be used, such as calling refresh manually from the view model itself when it thinks new data should arrive. For instance, after calling the run optimizer method, a refresh method is called for the Optimizer view model to populate fresh data. Such a method helped to solve the issue but due to its nature of not being so reactive to changes, many bugs have arisen where data was stale and caused desync. Adding more points to refresh data helped a bit more. In some cases, Data Manager was used to store data meant to be accessed between multiple components - such as current scenario and period. It worked but it caused additional issues with multiple scenarios and periods and required granular tweaks to make it be more reliable. In hindsight, paying the cost of changing many UI classes to pass around data to components could have been better and more reliable than relying more on Data Manager for passing data around between components. Furthermore, a clear and consistent way to pass data around Avalonia components from the project start could have resulted in better results. Even multiple attempts to refactor code were not able to resolve major issues of the codebase.\par
In addition to Avalonia C\# framework, LiveCharts library had its own issues too. After adding a few graphs with a total of around 3 thousand points, there were visible performance issues when resizing the application window. Because view model data did not change during application window resizing, it could have only been caused by the LiveCharts library. In the Optimizer page where there were more graphs than could fit in the screen, charts which were rendered off screen were hidden or distorted even when they scrolled into view. The hover tooltips would still show up but the visible chart had issues. The only workaround we found was resizing the window to refresh LiveCharts charts views. Furthermore, a LiveCharts feature of adding disjointed line sections was used. For example, when a unit is running, then is powered off and turned on again, the points should not be connected and the gap should be shown instead. However, the lines at the ends of intervals in the charts had higher values and caused confusion about what values are really at those endpoints. But the points were rendered correctly and the hover tooltip was showing correct values which could only mean there is a bug in the LiveCharts library itself.\par
Frontend was designed to be a balance between being user-friendly and having in-depth features. It caters to simple users by being graphically appealing, having simple summary statistics and summary graphs, limiting the number of pages visible and providing an intuitive interface. For advanced users, we have added a feature to create, edit or delete production units and to change scenario data dynamically. Furthermore, they can inspect the optimizer better with the help of detailed scenario graphs which show granular key metrics such as unit cost for 1 MWh, unit cost for unit and scheduling of each unit. We had two major Figma design iterations over the course of the project and a few of the changes were made directly in code after feedback from the stakeholders. Sprint reviews helped with improving the frontend incrementally by collecting stakeholder feedback at regular intervals and changing user interface based on the feedback.

% --------------------------------
\section{Backend Implementation}

Backend is composed from the Optimizer and Data Manager. Optimizer is responsible for calculating optimal cost for the provided time period and assets and Data Manager is responsible for exposing data layer to domain and presentation layers.\par
Optimizer is a central part of the backend. Data Manager configures optimizer by loading sources, assets and any scenario specific data. During Optimizer initialization, data sources and assets are validated. After validation, net cost for each unit at each hour is calculated. Net cost is the cost of 1 MWh of produced heat for a given unit at a specific hour. Net cost calculation takes into account production costs of the unit and any electricity produced or consumed by the unit during heat production. After initialization, the optimizer is started in a separate thread. Optimizer tries to find an optimal maintenance period for 1 production unit from the list of available units to be maintained. Optimizer iterates over all possible maintenance periods for each unit and selects the maintenance period with the lowest cost. For each possible maintenance period, the Optimizer step is run which calculates the cost, heat production, electricity consumption among other properties with the given maintenance period. During this step, maintenance period for the unit is applied and for each hour the units with the lowest cost are selected until the heat demand for an hour is filled. After heat demand for every hour is fulfilled, the result with granular detailed data and total summaries are returned back to the maintenance period iterator. The lowest cost maintenance period with its result data is stored in RDM. The main thread is notified about the completion of the optimization.\par
Currently, Data Manager belongs in the domain layer and it is responsible for exposing AM, SDM and RDM to optimizer and frontend. Data Manager does this by providing a thin wrapper over AM, SDM, RDM and by providing methods which abstract interacting with the optimizer and updating scenarios and periods. Data Manager implements singleton pattern using C\# static members because there is only one globally accessible instance of Data Manager.

% --------------------------------
\section{Data Management}

Explain how data was modelled and managed. Cover:

\begin{itemize}
    \item Data structure.
    \item Data validation and consistency considerations.
    \item How the data structure integrates with the backend
          (if relevant).
\end{itemize}